\section{Epistemic Gap}
\label{sec:epistemic}

The SKC model captures what an agent \textit{lacks}; the Epistemic Gap
captures what an agent \textit{thinks it has} versus what it actually has.
This distinction is critical for predicting hallucination, the most
dangerous failure mode in SOC environments.

\begin{definition}[Expressed Uncertainty]
  The expressed uncertainty $U(\A, \T) \in [0,1]$ is the complement of
  the confidence score the agent associates with its output:
  $U = 1 - \mathrm{conf}(\hat{y})$.
\end{definition}

\begin{definition}[Epistemic Gap]
  \begin{equation}
    \E(\A, \T) = \delta_K(\A, \T) \cdot (1 - U(\A, \T))
    \label{eq:epistemic_gap}
  \end{equation}
\end{definition}

Equation~\eqref{eq:epistemic_gap} formalises the \textit{confident ignorance}
failure mode. An agent with large knowledge deficit ($\delta_K \approx 1$)
that simultaneously reports high confidence ($U \approx 0$) produces
$\E \approx 1$: the maximal epistemic gap and the precondition for
hallucination. Conversely, an agent that correctly identifies its knowledge
limits ($U \approx 1$) produces $\E \approx 0$: it may fail the task but
will not fabricate an answer.

\begin{hypothesis}[Hallucination Precondition]
\label{hyp:hallucination}
  Hallucination probability $P_H$ is driven primarily by confident ignorance:
  $P_H \propto \E(\A, \T) = \delta_K \cdot (1 - U)$. This is empirically
  supported in Section~\ref{sec:experimental}.
\end{hypothesis}

\textbf{Connection to calibration theory.}
The Epistemic Gap operationalises findings from the confidence calibration
literature~\cite{guo2017calibration} for the specific case of agentic AI.
Modern neural networks are known to be systematically overconfident,
particularly on out-of-distribution inputs. The Epistemic Gap measures
whether this miscalibration occurs precisely where the agent's knowledge
is weakest: the most dangerous combination. A well-calibrated agent (one
whose confidence accurately reflects its accuracy) will produce low
$\E$ even under knowledge deficit, because it appropriately hedges its
outputs. This suggests that improving calibration in LLM-based security
agents may be as important as expanding their knowledge for reducing
dangerous failures.

\textbf{Why hallucination is particularly dangerous in SOC.}
An agent that says ``I don't have sufficient context'' is safe: it triggers
human escalation. An agent that confidently asserts a false threat
attribution is dangerous: it may trigger automated containment actions,
propagate fabricated IOCs to shared feeds, and erode analyst trust. The
Epistemic Gap captures exactly this distinction: it separates safe ignorance
(high $\delta_K$, high $U$, low $\E$) from dangerous confident ignorance
(high $\delta_K$, low $U$, high $\E$).

\textbf{Practical interpretation.}
Table~\ref{tab:eg_examples} illustrates the four quadrants of the Epistemic
Gap space with concrete SOC scenarios.

\begin{table}[t]
\caption{Epistemic Gap Quadrants with SOC Examples}
\label{tab:eg_examples}
\centering
\footnotesize
\setlength{\tabcolsep}{2.5pt}
\begin{tabular}{ccp{3.8cm}}
\toprule
$\delta_K$ & $U$ & \textbf{Scenario and Risk} \\
\midrule
Low  & Low  & Agent knows the threat and is confident. \textit{Safe: correct output likely.} \\
Low  & High & Agent knows but hedges. \textit{Suboptimal but safe: triggers review.} \\
High & High & Agent lacks knowledge and admits it. \textit{Safe: triggers escalation.} \\
High & Low  & Agent lacks knowledge but is confident. \textit{Dangerous: hallucination risk maximal ($\E \approx 1$).} \\
\bottomrule
\end{tabular}
\end{table}

The bottom-right quadrant (high $\delta_K$, low $U$, maximal $\E$) is
precisely where hallucination occurs: the agent does not possess the
knowledge required to answer correctly but proceeds with high confidence
regardless. The Epistemic Gap metric targets this quadrant specifically,
enabling monitoring systems to flag the most dangerous outputs before they
trigger automated actions.
